\documentclass[10pt,a4paper]{amsart}
\usepackage[margin=19mm]{geometry}
\usepackage{amsmath,amssymb,mathtools}
\usepackage[scr=rsfs,cal=euler]{mathalfa}
\usepackage{xfrac}
\usepackage[unicode,hidelinks,bookmarks=false]{hyperref}
\newcommand{\ie}{i.e.\ }
\newcommand{\cf}{cf.\ }
\newcommand{\resp}{respectively\ }
\newcommand{\Subsection}{\subsection}
\providecommand{\nobreakdash}{\nobreak\hbox{-}}
\setlength{\emergencystretch}{2em}
\multlinegap0pt
\usepackage{amssymb}
\usepackage{xcolor}
\usepackage{dsfont}
\newtheorem{lemma}{Lemma}
\newcommand{\N}{{\mathbb N}}
\newcommand{\R}{{\mathbb R}}
\renewcommand{\a}{\alpha}
\renewcommand{\d}{\delta}
\renewcommand{\div}{\mathrm{div}\,}
\DeclareMathOperator{\divergence}{div}
\newcommand{\e}{\varepsilon}
\DeclareMathOperator*{\esssup}{ess\,sup}
\newcommand{\grad}{\nabla}
\newcommand{\la}{\langle}
\newcommand{\lm}{\lambda}
\newcommand{\mA}{\mathcal{A}}
\newcommand{\pa}{\partial}
\newcommand{\ph}{\varphi}
\newcommand{\ra}{\rangle}
\title{On the Dirichlet-Neumann operator\\ for nearly spherical domains}
\author{Pietro Baldi, Vesa Julin, Domenico Angelo La Manna}
\date{Source: arXiv:2603.27865v1 (CC BY 4.0)}
\begin{document}
\maketitle

\noindent\emph{Source and licence.} On the Dirichlet-Neumann operator\\ for nearly spherical domains. Version arXiv:2603.27865v1 (CC BY 4.0), distributed by arXiv under the Creative Commons Attribution 4.0 International licence, http://creativecommons.org/licenses/by/4.0/, as recorded in the arXiv licence metadata for this exact version and shown on the abstract page https://arxiv.org/abs/2603.27865v1. Full licence notice: \texttt{licenses/arxiv-2603.27865-CC-BY-4.0.txt}.\par\smallskip

\noindent\textit{Editorial scope.} Counted unit: the article's lemma lemma:pinoli.S (source0.tex lines 1490--1511). The article's complete original proof of it is delivered with it (source0.tex lines 1513--1909, one \texttt{proof} environment of 397 lines). No mathematics is written, repaired or re-derived: the statement and the proof are the article's own lines, in the article's own notation, and no retained line is substituted or re-flowed.  Retained uncounted as the source-local context the proof needs: lines 687--689; lines 704--707; lines 807--812; lines 826--830; lines 966--972; lines 981--986; lines 990--999; lines 1060--1070; lines 1175--1179; lines 1182--1185; lines 1189--1193; lines 1213--1218; lines 1229--1244; lines 1363--1375; lines 1401--1406; lines 1418--1421; lines 1440--1444; lines 1456--1459; lines 1467--1470.  Nothing was omitted from any retained span, so the package carries no omission record.  No retained line contains a citation, so the wrapper carries no bibliography and no bibliography entry.  No retained line contains a figure environment, an image-inclusion command or drawing code, so no image is delivered and the package declares no asset dependency.  Editorial formatting only, in addition: an AMS \texttt{amsart} preamble that restates the article's own declarations and macros used by the retained lines, namely the environments \texttt{lemma} on separate counters, together with the standard packages those lines need.  The retained environments print in this document as Lemma 1 in order of appearance, whereas the article prints the same items with the numbering of its own full document.  The complete native source of the article as distributed in the arXiv e-print for this exact version is bundled unchanged as \texttt{sources/197391/source0.tex}.
\begin{equation} \label{res.unity.in.the.annulus}
\sum_{j=1}^N \psi_j(x) = 1 \quad \forall x \in \R^n, \  1-\delta \leq |x| \leq 1.
\end{equation}

\begin{equation} \label{explicit.mathtt.f.inv}
\mathtt g : \mathbb H_1 \to \mathbb H_0, \quad \ 
\mathtt g(y) := (1 - y_n) (1 + |y'|^2)^{-\frac12} (y', -1),
\end{equation}

\begin{equation}  \label{formula.transition.map}
\mathtt f_\ell \circ \mathtt g_j(y) = (\mathtt T_{\ell j}(y') , y_n), 
\quad \ 
\mathtt T_{\ell j}(y') 
= - \frac{[\mathtt R_\ell \mathtt R_j^{-1} (y', -1)]'}{[\mathtt R_\ell \mathtt R_j^{-1} (y', -1)]_n}
\end{equation}

\begin{equation} \label{def.tilde.mathtt.f}
\tilde{\mathtt f}(x) := 
x - p_0 + \ph_\delta(x-p_0) ( \mathtt f(x) - x + p_0 )
\quad \ \forall x \in \R^n,
\end{equation}

\begin{equation}  \label{rotat.09}
\tilde{\mathtt T}_{\ell j}(y') 
:= - \lm \mathtt b 
+ \lm^2 (\mathtt A \mathtt d - \mathtt b \mathtt c ) y' \big( 1 + \ph_\d(y') S(y') \big), 
\quad \ 
\ph_\d(y') := \ph \Big( \frac{|y'|}{4 \delta} \Big),
\end{equation}

\begin{equation} \label{extension.in.A0}
(\tilde{\mathtt T}_{\ell j}(y'), y_n) 
= (\mathtt T_{\ell j}(y'), y_n) 
= \mathtt f_\ell (\mathtt g_j(y)) 
\quad \ \forall y = (y', y_n) \in A_0.
\end{equation} 

\begin{lemma}  \label{lemma:rotat.tilde}
There exists a universal constant $\delta_0 \in (0, 1/8)$ such that, 
for $0 < \delta \leq \delta_0$, if $K_j \cap K_\ell \neq \emptyset$, then 
the map $\tilde{\mathtt T}_{\ell j}$ in \eqref{rotat.09}
is a $C^\infty$ diffeomorphism of $\R^{n-1}$ 
that coincides with $\mathtt T_{\ell j}$ in \eqref{formula.transition.map} 
in the ball $|y'| \leq 4 \delta$,
and coincides in $|y'| \geq 8 \delta$ with the affine map 
$y' \mapsto - \lm \mathtt b + \lm^2 (\mathtt A \mathtt d - \mathtt b \mathtt c ) y'$.
\end{lemma}

\begin{align} 
\| \pa_{x_k} u \|_{H^{s,r}(\R^n_+)} 
& \leq \| u \|_{H^{s,r+1}(\R^n_+)} 
\leq \| u \|_{H^{s+1,r}(\R^n_+)}, 
\quad k=1, \ldots, n-1,
\label{pa.k.est}
\\
\| \pa_{x_n} u \|_{H^{s,r}(\R^n_+)} 
& \leq \| u \|_{H^{s+1,r}(\R^n_+)},
\label{pa.n.est}
\end{align} 

\begin{equation} \label{est.prod.infty.0}
\| fu \|_{H^{0,r}(\R^n_+)} 
\leq 2 \| f \|_{L^\infty(\R^n_+)} \| u \|_{H^{0,r}(\R^n_+)} 
+ C_r \| f \|_{W^{0,b,\infty}(\R^n_+)} \| u \|_{H^{0,0}(\R^n_+)}
\end{equation}

\begin{equation}  \label{def.norm.W.0.b.infty}
\| f \|_{W^{0,b,\infty}(\R^n_+)} := 
\esssup_{x_n \in (0,\infty)} \| f(\cdot, x_n) \|_{W^{b,\infty}(\R^{n-1})}, 
\end{equation}

\begin{equation} \label{est.prod.infty.1}
\| fu \|_{H^{1,r}(\R^n_+)} 
\leq C_0 \| f \|_{W^{1,\infty}(\R^n_+)} \| u \|_{H^{1,r}(\R^n_+)} 
+ C_r \| f \|_{W^{1,b,\infty}(\R^n_+)} \| u \|_{H^{1,0}(\R^n_+)}
\end{equation}

\begin{equation} 
\| x \mapsto u (f(x'),x_n) \|_{H^{s,r}(\R^n_+)} 
\leq C_{r,f} ( \| u \|_{H^{s,r}(\R^n_+)} 
+ \| g \|_{H^{1 + r_0 + r + s}(\R^{n-1})} \| u \|_{H^{s,0}(\R^n_+)} ) 
\label{composition.est.s.r} 
\end{equation}

\begin{lemma} \label{lemma:est.Lap.R.n.+}
Let $r \geq 0$ be any real number, let $s=0$ or $s=1$, 
and let $u \in H^{s+1,r}(\R^n_+)$. Then 
\begin{align} 
\| u \|_{H^{s+1,r}(\R^n_+)} 
& \leq C \| \Delta u \|_{H^{s-1,r}(\R^n_+)} 
+ C^{r+1} \| \Delta u \|_{H^{s-1,0}(\R^n_+)} 
+ C \| u(\cdot, 0) \|_{H^{s + \frac12 + r}(\R^{n-1})}
\notag \\ 
& \quad \ 
+ C^{r+1} \| u(\cdot, 0) \|_{H^{s + \frac12}(\R^{n-1})}
+ C^{r+1} \| u \|_{H^{s,0}(\R^n_+)} 
\label{meran.04}
\end{align}
where the constant $C$ is independent of $s, r, u$. 
\end{lemma}

\begin{lemma}[Solution map for the Poisson problem in divergence form] 
\label{lemma:op.S}
For any $g \in L^2(B_1)$ there exists a unique $u$ such that 
\[
u \in H^1_0(B_1), \quad \ \  
\Delta u = \div g \ \  \text{in } B_1
\]
in the sense of distributions. 
The map $g \to u$ is a linear continuous operator 
\[
\mathtt S : L^2(B_1) \to H^1_0(B_1).
\]
\end{lemma}

\begin{equation}  \label{def.rho.k.zeta.k}
\rho_k := 1 - \frac{\delta}{4} \sum_{\ell=0}^k \frac{1}{2^\ell}, 
\quad 
\zeta_k(x) := \ph \Big( \frac{|x| - \rho_{k+1}}{\rho_k - \rho_{k+1}} \Big)
\quad \forall x \in \R^n, \ \ k \in \N_0.
\end{equation}

\begin{equation} \label{zeta.k.is.1.where}
\zeta_{k+1} = 1 \text{ on } \mathrm{supp}(\zeta_k), 
\quad \ \forall k \in \N_0.
\end{equation}

\begin{align} \label{def.seminorm}
|u|_{X^{s,r}_k} := \sum_{j=1}^N \| (u \zeta_k \psi_j) \circ \mathtt g_j \|_{H^{s,r}(\R^n_+)},
\quad 
|u|_{X^{s,r}_{*,k}} := \sum_{j=1}^N \| (u \zeta^*_k \psi_j) \circ \mathtt g_j \|_{H^{s,r}(\R^n_+)},
\end{align}

\begin{align} \label{est.u.X.low}
|u|_{X^{0,0}_k} \leq C \| u \|_{L^2(B_1)}, \quad \ 
|u|_{X^{1,0}_k} \leq C_k \| u \|_{H^1(B_1)}, 
\end{align}

\begin{equation}  \label{est.zeta.k.psi.j}
\| (\zeta_k \psi_j) \circ \mathtt g_j \|_{L^\infty(\R^n_+)} \leq 1, \quad 
\| \pa_y^\a \{ (\zeta_k \psi_j) \circ \mathtt g_j \} \|_{L^\infty(\R^n_+)} \leq C_\a 2^{k|\a|}. 
\end{equation}

\begin{lemma} \label{lemma:pinoli.S}
The linear map $\mathtt S$ defined in Lemma \ref{lemma:op.S} satisfies 
\begin{align}
|\mathtt S g|_{X^{2,0}_k} 
& \leq C_0 |g|_{X^{1,0}_k} + C_k \| g \|_{L^2(B_1)} 
\label{pinoli.15}
\end{align}
and, for all integers $b \geq 1$, all real $r$ in the interval $b-1 < r \leq b$, 
\begin{align} 
|\mathtt S g|_{X^{2,r}_k} \leq 
C_0 |g|_{X^{1,r}_k}
+ C_{r,k} \Big( \| g \|_{L^2(B_1)}
+ |g|_{X^{1,0}_k}
+ |g|_{X^{1,0}_{k+b}}
+ \sum_{\ell=1}^{b-1} |g|_{X^{1,r-\ell}_{k+\ell}} \Big),
\label{pinoli.16}
\end{align}
where $C_0$ is independent of $b,r,k$,  
and $C_{r,k}$ is increasing in $r$ and in $k$. 
For $b=1$ the last sum is empty.
The same also holds for the $*$ seminorms. 
\end{lemma}

\begin{proof}
Let $u = \mathtt S g$,   
so that $u \in H^1_0(B_1)$ and $\Delta u = \div g$ in $B_1$. 
Let $j \in \{1, \ldots, N\}$, and consider the diffeomorphism 
$\mathtt f_j : B_1 \cap K_j \to A_0 \cap \R^n_+$ 
and its inverse $\mathtt g_j$. 
Since $\Delta u = \div g$, 
taking test functions in $C^{\infty}_c(B_1 \cap K_j)$, 
integrating in $B_1$, using the divergence theorem, 
making the change of variable $x = \mathtt g_j(y)$ in the integral, 
and using the transformation rule for the gradient
\begin{equation} \label{trans.rule.grad}
(\grad u) \circ \mathtt g_j (y) = [D \mathtt g_j(y)]^{-T} \grad( u \circ \mathtt g_j)(y), 
\quad \ y \in A_0 \cap \R^n_+,
\end{equation}
one finds that 
\begin{equation} \label{eq.for.u.circ.mathtt.g.j}
\divergence (\mathtt p \mathtt Q \grad \tilde u_j) 
= \divergence (\mathtt p (D \mathtt g_j)^{-1} \tilde g_j) 
\quad \text{in } A_0 \cap \R^n_+,
\end{equation}
where 
\begin{alignat}{2} 
\tilde u_j & := u \circ \mathtt g_j, \qquad & 
\mathtt p(y) & := \det (D \mathtt g_j(y)) = \det (D \mathtt g(y)), 
\label{def.mathtt.p}
\\
\tilde g_j & := g \circ \mathtt g_j, \qquad & 
\mathtt Q(y) & := [D \mathtt g_j(y)]^{-1} [D \mathtt g_j(y)]^{-T} 
= [D \mathtt g(y)]^{-1} [D \mathtt g(y)]^{-T}
\label{def.mathtt.Q}
\end{alignat}
in $A_0 \cap \R^n_+$. 
The last identities in \eqref{def.mathtt.p} and \eqref{def.mathtt.Q}
hold because $\mathtt g_j = \mathtt R_j^{-1} \circ \mathtt g$ 
and $\mathtt R_j$ is a rotation. 
By \eqref{explicit.mathtt.f.inv}, in $A_0$ one has 
\begin{align}
\label{una.buona.volta}
\mathtt Q(y) & = \begin{bmatrix} 
\mathtt q(y) (\mathrm{I'} + y' \otimes y')  & 0 \\
0 & 1 \end{bmatrix}\!\!,
\quad
\mathtt q(y) := \frac{1+|y'|^2}{(1-y_n)^2},
\quad  
\mathtt p(y) = \frac{ (1-y_n)^{n-1} }{ (1 + |y'|^2)^{\frac{n}{2}} },
\end{align}
where $I'$ is the $(n-1) \times (n-1)$ identity matrix. 
Since the bottom-right entry of the matrix $\mathtt Q$ is 1, 
and since derivatives with respect to $y_n$ play a separate role, 
it is convenient to work with $\mathtt Q$ instead of $\mathtt P$. 
Dividing by $\mathtt p$, \eqref{eq.for.u.circ.mathtt.g.j} becomes 
\begin{align} 
\divergence (\mathtt Q \grad \tilde u_j) 
& = - \mathtt p^{-1} \la \grad \mathtt p , \mathtt Q \grad \tilde u_j \ra
+ \mathtt p^{-1} \divergence (\mathtt p (D \mathtt g_j)^{-1} \tilde g_j) 
\quad \text{in } A_0 \cap \R^3_+.
\label{eq.for.u.circ.mathtt.g.j.with.mathtt.Q}
\end{align}
Now, given any two functions $v,a$, one has 
\begin{equation} \label{div.P.grad.v.m}
\divergence (\mathtt Q \grad (va)) 
= a \divergence (\mathtt Q \grad v) 
+ 2 \la \mathtt Q \grad a , \grad v \ra 
+ v \divergence (\mathtt Q \grad a).
\end{equation}
To take advantage of the fact that $\mathtt Q(y)$ is close to $\mathtt Q(0) = I$ for $y$ close to $0$, 
we write $\mathtt Q$ in the LHS of \eqref{div.P.grad.v.m} 
as the sum $I + (\mathtt Q - I)$, and \eqref{div.P.grad.v.m} becomes 
\begin{equation} \label{Lap.v.m}
\Delta (va)
= \divergence \{ (I - \mathtt Q) \grad (va) \}
+ a \divergence (\mathtt Q \grad v) 
+ 2 \la \mathtt Q \grad a , \grad v \ra 
+ v \divergence ( \mathtt Q \grad a).
\end{equation}
Let 
\begin{alignat}{2}
\tilde \psi_j & := \psi_j \circ \mathtt g_j, \qquad & 
a_{kj} & := \tilde \zeta_{kj} \tilde \psi_j = (\zeta_k \psi_j) \circ \mathtt g_j, 
\label{def.a.kj}
\\ 
\tilde \zeta_{kj} & := \zeta_k \circ \mathtt g_j, \qquad &  
w_{kj} & := \tilde u_j a_{kj} = \tilde u_j \tilde \zeta_{kj} \tilde \psi_j
= (u \zeta_k \psi_j) \circ \mathtt g_j.
\label{def.w.kj}
\end{alignat} 
By \eqref{Lap.v.m} with $v = \tilde u_j$, $a = a_{kj}$, 
using \eqref{eq.for.u.circ.mathtt.g.j.with.mathtt.Q} 
to substitute the second term in the RHS of \eqref{Lap.v.m}, 
we get 
\begin{align} 
\Delta w_{kj} 
& = \divergence ((I - \mathtt Q) \grad w_{kj} ) 
- a_{kj} \mathtt p^{-1} \la \grad \mathtt p , \mathtt Q \grad \tilde u_j \ra
+ a_{kj} \mathtt p^{-1} \divergence (\mathtt p (D \mathtt g_j)^{-1} \tilde g_j) 
\notag \\ & \quad  
+ 2 \la \mathtt Q \grad a_{kj} , \grad \tilde u_j \ra
+ \tilde u_j \divergence ( \mathtt Q \grad a_{kj} )
\label{eq.Lap.w}
\end{align}
in $A_0 \cap \R^n_+$. 
Note that $\tilde \psi_j$ is compactly supported in $A_0$ 
because $\psi_j$ is compactly supported in $K_j$. 
Hence $w_{kj}$ is also compactly supported in $A_0$, and we 
extend it to $\R^n_+$ by setting $w_j := 0$ in $\R^n_+ \setminus A_0$. 
We also extend $\mathtt Q, \mathtt p$ to $\R^n$ 
replacing $\mathtt g$ with its extension 
$\tilde{\mathtt g} = \tilde{\mathtt f}^{-1}$, see \eqref{def.tilde.mathtt.f};
with a little abuse of notation, we denote the extensions without changing letters.
For future convenience, we note here that, 
by \eqref{def.tilde.mathtt.f}, one has $\tilde{\mathtt f}(x) = x - p_0$ 
for all $x$ outside the ball $B_{4\delta}(p_0)$, 
and therefore $\mathtt p = 1$, $\mathtt Q = I$ 
outside the set $\mathtt f(B_{4\delta}(p_0))$, 
which is a neighborhood of $A_0 = \mathtt f(B_{2\delta}(p_0))$. 

Since $u \in H^1_0(B_1)$ and $\psi_j, \zeta_k, \mathtt g_j$ are smooth, 
one has $w_{kj} \in H^1_0(\R^n_+)$. 
Then, by Lemma \ref{lemma:est.Lap.R.n.+}, for every real $r \geq 0$ one has 
\begin{align} 
\| w_{kj} \|_{H^{2,r}(\R^n_+)} 
& \leq C \| \Delta w_{kj} \|_{H^{0,r}(\R^n_+)} 
+ C^{r+1} \| \Delta w_{kj} \|_{L^2(\R^n_+)} 
+ C^{r+1} \| w_{kj} \|_{H^1(\R^n_+)}, 
\label{olive.01}
\end{align}
where $C > 0$ is the constant in \eqref{meran.04},
and $\Delta w_{kj}$ is given by \eqref{eq.Lap.w}. 
Now we estimate each term in the right-hand side of \eqref{eq.Lap.w}. 

\emph{First term}. 
By \eqref{una.buona.volta}, the last column and the last row 
of the matrix $(I - \mathtt Q(y))$ vanish in $A_0$ $\cap \, \R^n_+$,
and $\divergence ((I - \mathtt Q) \grad w_j)$ 
is the sum of terms of the form $\pa_{y_i} ( c(y) \pa_{y_k} w_j)$ 
with $i, k \in \{ 1, \ldots, n-1 \}$, without derivatives of $w_j$ with respect to $y_n$. 
Hence, by \eqref{pa.k.est} and \eqref{est.prod.infty.0}, one has 
\begin{align*}
\| \divergence ((I - \mathtt Q) \grad w_{kj}) \|_{H^{0,r}(\R^n_+)} 
& \leq C_0 \| I - \mathtt Q \|_{L^\infty(\R^n_+)} \| w_{kj} \|_{H^{0,r+2}(\R^n_+)}
\notag \\ & \quad \ 
+ C_r \| I - \mathtt Q \|_{W^{0,b,\infty}(\R^n_+)} \| w_{kj} \|_{H^{0,1}(\R^n_+)},
\end{align*}
where $C_0, C_r, b$ are given by the application of \eqref{est.prod.infty.0}. 
Given any $\e > 0$, the radius $2\delta$ of the balls $K_1, \ldots, K_N$ 
can be chosen sufficiently small to have 
$C_0 \| I - \mathtt Q \|_{L^\infty(\R^n_+)} \leq \e$; 
note that such a $\delta$ is independent on $r$. 
Hence 
\begin{align}  \label{olive.02}
\| \divergence ((I - \mathtt Q) \grad w_{kj}) \|_{H^{0,r}(\R^n_+)} 
& \leq \e \| w_{kj} \|_{H^{0,r+2}(\R^n_+)} 
+ C_r \| w_{kj} \|_{H^{0,1}(\R^n_+)},
\end{align} 
where $C_r$ is a constant depending on $r$, increasing in $r$, 
incorporating the factor $\| I - \mathtt Q \|_{W^{0,b,\infty}(\R^n_+)}$. 

\emph{Second term}. 
By \eqref{def.mathtt.p}, 
\eqref{zeta.k.is.1.where} and \eqref{res.unity.in.the.annulus}, 
the second term in the RHS of \eqref{eq.Lap.w} is 
\begin{align}
& a_{kj} \mathtt p^{-1} 
\la \grad \mathtt p , \mathtt Q \grad \tilde u_{j} \ra
= \sum_{\ell \in \mA_j} a_{kj} \mathtt p^{-1} \la \grad \mathtt p , \mathtt Q \grad 
\{ (u \zeta_{k+1} \psi_\ell) \circ \mathtt g_j \} \ra,
\label{echin.03}
\end{align}
where $\mA_j := \{ \ell \in \{ 1, \ldots, N \} : K_\ell \cap K_j \neq \emptyset \}$. 
Identity \eqref{echin.03} holds because, by \eqref{zeta.k.is.1.where}, 
$\zeta_{k+1} \circ \mathtt g_j = 1$ in the support of $\zeta_k \circ \mathtt g_j$, 
by \eqref{res.unity.in.the.annulus},  
and because, if the balls $K_\ell, K_j$ are disjoint, 
then $\psi_\ell \circ \mathtt g_j = 0$ in the support of $\psi_j \circ \mathtt g_j$. 
Moreover, for every $\ell \in \mA_j$, one has 
\begin{equation} \label{transition.w}
(u \zeta_{k+1} \psi_\ell) \circ \mathtt g_j 
= (u \zeta_{k+1} \psi_\ell) \circ \mathtt g_\ell \circ \mathtt f_\ell \circ \mathtt g_j
= w_{k+1,\ell} \circ (\mathtt f_\ell \circ \mathtt g_j),
\end{equation}
where the transition map $\mathtt f_\ell \circ \mathtt g_j$ 
in $A_0$ is given by formula \eqref{formula.transition.map}.
We consider the extension of the transition map to $\R^n_+$ 
given in \eqref{rotat.09}, \eqref{extension.in.A0}. 
By Lemma \ref{lemma:rotat.tilde}, the function $\tilde{\mathtt T}_{\ell j}$ 
is a diffeomorphism of $\R^{n-1}$, and it is the sum 
of an affine map and a smooth function with compact support. 
Hence, by the composition estimate \eqref{composition.est.s.r}, we get
\begin{equation}  \label{echin.04}
\| w_{k+1,\ell} \circ (\mathtt f_\ell \circ \mathtt g_j) \|_{H^{1,r}(\R^n_+)}
\leq C_r \|  w_{k+1,\ell} \|_{H^{1,r}(\R^n_+)},
\end{equation}
where the constant $C_r$ is increasing in $r$. 
Concerning $a_{kj}$, one has estimate \eqref{est.zeta.k.psi.j}. 
Also, by \eqref{def.rho.k.zeta.k}, 
$\zeta_k(\mathtt g_j(y))$ depends only on $|\mathtt g_j(y)|$; 
moreover $\mathtt g_j(y) = \mathtt R_j^{-1} \mathtt g(y)$, 
and $|\mathtt g_j(y)| = |\mathtt g(y)| = 1 - y_n$ by \eqref{explicit.mathtt.f.inv}.  
Hence $\zeta_k \circ \mathtt g_j$ depends only on $y_n$, and not on $y'$, 
so that $\pa_y^\a \{ (\zeta_k \psi_j) \circ \mathtt g_j \} 
= (\zeta_k \circ \mathtt g_j) \pa_y^\a (\psi_j \circ \mathtt g_j)$ 
for all multi-indices $\a = (\a', 0)$. 
As a consequence, by \eqref{def.a.kj} and \eqref{def.norm.W.0.b.infty}, one has
\begin{equation}  \label{est.a.W.0.b.infty}
\| a_{kj} \|_{W^{0,b,\infty}(\R^n_+)}
= \| (\zeta_k \psi_j) \circ \mathtt g_j \|_{W^{0,b,\infty}(\R^n_+)}
\leq \| \psi_j \circ \mathtt g_j \|_{W^{0,b,\infty}(\R^n_+)} \leq C_r
\end{equation}
for some constant $C_r$ 
dependent on $r$, increasing in $r$ (in fact, $C_r$ depends on $b$).  
By \eqref{est.prod.infty.0}, \eqref{transition.w}, \eqref{echin.04}, 
\eqref{est.zeta.k.psi.j}, \eqref{est.a.W.0.b.infty}, 
each term of the sum in the RHS of \eqref{echin.03} satisfies 
\begin{equation}  \label{echin.05}
\| a_{kj} \mathtt p^{-1} \la \grad \mathtt p , \mathtt Q \grad 
\{ (u \zeta_{k+1} \psi_\ell) \circ \mathtt g_j \} \ra \|_{H^{0,r}(\R^n_+)}
\leq C_r \| w_{k+1,\ell} \|_{H^{1,r}(\R^n_+)}.
\end{equation}
The sum of \eqref{echin.05} over $\ell \in \mA_j$ gives an estimate 
of the second term in the RHS of \eqref{eq.Lap.w}.

\emph{Third term.} We write the third term in the RHS of \eqref{eq.Lap.w} as 
\begin{align}
a_{kj} \mathtt p^{-1} \divergence (\mathtt p (D \mathtt g_j)^{-1} \tilde g_j) 
= \mathtt p^{-1} \divergence (\mathtt p (D \mathtt g_j)^{-1} \gamma_{kj}) 
- \la (D \mathtt g_j)^{-1} \tilde g_j , \grad a_{kj} \ra,
\label{third.term.id}
\end{align}
where 
\begin{equation} \label{def.gamma.nkj}
\gamma_{kj} := \tilde g_j a_{kj} 
= (g \zeta_k \psi_j) \circ \mathtt g_j.
\end{equation}
By \eqref{est.prod.infty.0}, \eqref{pa.k.est}, \eqref{pa.n.est}, 
\eqref{est.prod.infty.1}, the first term in the RHS of \eqref{third.term.id} satisfies 
\begin{equation}  \label{olive.21}
\| \mathtt p^{-1} \divergence (\mathtt p (D \mathtt g_j)^{-1} \gamma_{kj}) \|_{H^{0,r}(\R^n_+)} 
\leq C_0 \| \gamma_{kj} \|_{H^{1,r}(\R^n_+)} + C_r \| \gamma_{kj} \|_{H^{1,0}(\R^n_+)}.
\end{equation}
Using \eqref{def.mathtt.Q}, \eqref{zeta.k.is.1.where} and \eqref{res.unity.in.the.annulus} 
like in \eqref{echin.03}, and using the identity 
$\mathtt g_j = \mathtt g_\ell \circ \mathtt f_\ell \circ \mathtt g_j$ as in \eqref{transition.w},
we write the last term of \eqref{third.term.id} as 
\begin{align}
\la (D \mathtt g_j)^{-1} \tilde g_j , \grad a_{kj} \ra
& = \sum_{\ell \in \mA_j} 
\la (D \mathtt g_j)^{-1} ((g \zeta_{k+1} \psi_\ell) \circ \mathtt g_j) , \grad a_{kj} \ra
\notag \\
& = \sum_{\ell \in \mA_j} 
\la (D \mathtt g_j)^{-1} (\gamma_{k+1, \ell} \circ (\mathtt f_\ell \circ \mathtt g_j)) , \grad a_{kj} \ra.
\label{echin.06}
\end{align}
Hence, using \eqref{est.prod.infty.0}, 
\eqref{est.zeta.k.psi.j}, 
\eqref{est.a.W.0.b.infty}, 
and Lemma \ref{lemma:rotat.tilde} and \eqref{composition.est.s.r} 
for the composition with the transition map $\mathtt f_\ell \circ \mathtt g_j$,
one proves that each term of the sum in \eqref{echin.06} satisfies 
\begin{equation} \label{olive.20}
\| \la (D \mathtt g_j)^{-1} (\gamma_{k+1, \ell} \circ (\mathtt f_\ell \circ \mathtt g_j)) , 
\grad a_{kj} \ra \|_{H^{0,r}(\R^n_+)} 
\leq C_r 2^k \| \gamma_{k+1, \ell} \|_{H^{0,r}(\R^n_+)}.
\end{equation}
The constant $C_r$ in \eqref{olive.20} depends on $r$ 
because we have estimated a composition; 
the factor $2^k$ is due to the presence 
of the first derivative $\pa_{y_n} (\zeta_{k+1} \circ \mathtt g_j)$.  
From \eqref{olive.21} and \eqref{olive.20} we obtain an estimate for \eqref{third.term.id}.


\emph{Fourth and fifth term}. 
Proceeding as in \eqref{echin.03}-\eqref{echin.05}, one proves that the 
fourth and fifth term in \eqref{eq.Lap.w} are the sum over $\ell \in \mA_j$ of terms 
satisfying the same bound \eqref{echin.05}, but with $C_r 2^{2k}$ instead of $C_r$. 
The factor $2^{2k}$ is due to the presence of the second derivative of $\zeta_{k+1}$ 
with respect to $y_n$. 
 

\emph{Estimate of $u$}. 
By \eqref{eq.Lap.w}, \eqref{olive.02}, \eqref{echin.05}, \eqref{olive.21}, \eqref{olive.20}, 
and \eqref{echin.05} with factor $2^{2k}$, one has 
\begin{multline}
\| \Delta w_{kj} \|_{H^{0,r}(\R^n_+)}
\leq 
\e \| w_{kj} \|_{H^{0, r+2}(\R^n_+)}
+ C_r \| w_{kj} \|_{H^{1,0}(\R^n_+)}
+ C_r 2^{2k} \sum_{\ell \in \mA_j} \| w_{k+1, \ell} \|_{H^{1,r}(\R^n_+)}
\\
+ C_0 \| \gamma_{kj} \|_{H^{1,r}(\R^n_+)} 
+ C_r \| \gamma_{kj} \|_{H^{1,0}(\R^n_+)} 
+ C_r 2^k \sum_{\ell \in \mA_j} \| \gamma_{k+1,\ell} \|_{H^{0,r}(\R^n_+)}.
\label{echin.10}
\end{multline}
By \eqref{olive.01}, \eqref{echin.10} and \eqref{echin.10}$|_{r=0}$, we get
\begin{align}
\| w_{kj} \|_{H^{2,r}(\R^n_+)}
& \leq 
\e C_0 \| w_{kj} \|_{H^{0, r+2}(\R^n_+)} 
+ \e C_r \| w_{kj} \|_{H^{0, 2}(\R^n_+)}
+ C_r \| w_{kj} \|_{H^{1,0}(\R^n_+)}
\notag \\ & \quad \ 
+ C_r 2^{2k} \sum_{\ell \in \mA_j} \| w_{k+1, \ell} \|_{H^{1,r}(\R^n_+)}
+ C_0 \| \gamma_{kj} \|_{H^{1,r}(\R^n_+)} 
+ C_r \| \gamma_{kj} \|_{H^{1,0}(\R^n_+)} 
\notag \\ & \quad \ 
+ C_r 2^k \sum_{\ell \in \mA_j} \| \gamma_{k+1,\ell} \|_{H^{0,r}(\R^n_+)}
\label{olive.05}
\end{align}
for some $C_0, C_r$ (possibly different than above).
Recalling definitions \eqref{def.w.kj}, \eqref{def.gamma.nkj} of $w_{kj}, \gamma_{kj}$,  
taking the sum of \eqref{olive.05} over $j = 1, \ldots, N$, 
and using notation \eqref{def.seminorm}, 
we obtain 
\begin{align}
|u|_{X^{2,r}_k}
& \leq \e a_0 |u|_{X^{0,r+2}_k}
+ \e b_r |u|_{X^{0,2}_k}
+ C_r |u|_{X^{1,0}_k}
+ C_r 2^{2k} |u|_{X^{1,r}_{k+1}}
\notag \\ & \quad \ 
+ C_0 |g|_{X^{1,r}_k}
+ C_r |g|_{X^{1,0}_k}
+ C_r 2^k |g|_{X^{0,r}_{k+1}},
\label{olive.09}
\end{align}
where we have denoted by $a_0$ and $b_r$ the constants appearing with the factor $\e$;
also, $b_r, C_r$ are increasing functions of $r$. 
For $r=0$, \eqref{olive.09} becomes
\begin{align}
|u|_{X^{2,0}_k}
& \leq \e (a_0 + b_0) |u|_{X^{0,2}_k}
+ C_0 |u|_{X^{1,0}_k}
+ C_0 2^{2k} |u|_{X^{1,0}_{k+1}}
+ 2 C_0 |g|_{X^{1,0}_k}
+ C_0 2^k |g|_{X^{0,0}_{k+1}}.
\label{olive.10}
\end{align}
Since $|u|_{X^{0,2}_k} \leq |u|_{X^{2,0}_k}$, 
we fix $\e$ (and therefore the half-radius $\delta$ of the balls $K_j$ 
and their cardinality $N$) 
such that $\e (a_0 + b_0) \leq 1/2$,
so that the first term in the RHS of \eqref{olive.10} can be absorbed 
by the LHS of \eqref{olive.10}. 
The second, the third and the fifth term in the RHS of \eqref{olive.10} are estimated 
by \eqref{est.u.X.low} and Lemma \ref{lemma:op.S}, and we get \eqref{pinoli.15}. 
Note that $\e, \delta, N$ do not depend on $r,j,k$. 

Now that $\e$ has been fixed, consider \eqref{olive.09} again. 
Since $a_0 \e \leq 1/2$ and $|u|_{X^{0,r+2}_k} \leq |u|_{X^{2,r}_k}$, 
the first term in the RHS of \eqref{olive.09} can be absorbed by the LHS of \eqref{olive.09}.  
Also, we use \eqref{pinoli.15} to estimate the second and third term in the RHS of \eqref{olive.09}. 
Thus, for all $r,k$, we get
\begin{align}
|u|_{X^{2,r}_k}
& \leq 
C_0 |g|_{X^{1,r}_k} + C_r |g|_{X^{1,0}_k} 
+ C_{r,k} (\| g \|_{L^2(B_1)} + |g|_{X^{0,r}_{k+1}} + |u|_{X^{1,r}_{k+1}} ).
\label{pinoli.12}
\end{align}
For $0 < r \leq 1$, we consider \eqref{pinoli.12}, 
we use the fact that 
\[
|g|_{X^{0,r}_{k+1}} \leq 
|g|_{X^{0,1}_{k+1}} \leq 
|g|_{X^{1,0}_{k+1}}, 
\quad 
|u|_{X^{1,r}_{k+1}} 
\leq |u|_{X^{1,1}_{k+1}}
\leq |u|_{X^{2,0}_{k+1}},
\]
and we use \eqref{pinoli.15} to bound $|u|_{X^{2,0}_{k+1}}$; 
in this way we obtain 
\begin{align}
|u|_{X^{2,r}_k} 
& \leq 
C_0 |g|_{X^{1,r}_k}
+ C_{r,k} (\| g \|_{L^2(B_1)} + |g|_{X^{1,0}_k} + |g|_{X^{1,0}_{k+1}}), 
\quad 0 < r \leq 1,
\label{pinoli.13}
\end{align}
with some possibly larger $C_{r,k}$, and the same $C_0$ of \eqref{pinoli.12}. 

Now we prove \eqref{pinoli.16} by induction on $b$. 
For $b=1$, \eqref{pinoli.16} is \eqref{pinoli.13}. 
Assume that \eqref{pinoli.16} is proved for some integer $b \geq 1$, 
and let $b < r \leq b+1$. 
Apply \eqref{pinoli.12}, 
use the inequality $|u|_{X^{1,r}_{k+1}} \leq |u|_{X^{2,r-1}_{k+1}}$
and the induction assumption \eqref{pinoli.16} (with $r-1, k+1$ instead of $r,k$)
to estimate $|u|_{X^{2,r-1}_{k+1}}$. This gives \eqref{pinoli.16} with $b+1$ instead of $b$ 
(and with constants $C_{r,k}$ possibly different than those of the induction assumption. 
This is not a problem, because now $r$ is in the next interval $(b, b+1]$, 
and each of such intervals has its own set of constants). 
Note that the constant $C_0$ in \eqref{pinoli.16} is not changed by the induction step. 
Hence \eqref{pinoli.16} holds for all integers $b \geq 1$. 
\end{proof}

\end{document}
